% q0029.tex — Music or Noise?
\question{
  id        = {0029},
  title     = {Music or Noise?},
  source    = {OBECON},
  year      = {2026},
  round     = {Seletiva IEO -- Phase B},
  language  = {English},
  topic     = {Micro},
  subtopic  = {Externalities / Social Welfare},
  type      = {dissertative},
  statement = {%
    Lucas likes to play music, but noise bothers his community --- even
    though they appreciate some amount of music. Every day, he chooses a
    proportion of the day to spend playing music, keeping the rest silent.
    His utility function is
    \[
      u(s,m) = 1 - (s^2 + m^2),
    \]
    where $s \in [0,1]$ is the proportion of the day spent in silence and
    $m \in [0,1]$ is the proportion spent playing music. Because time is
    split between silence and music:
    \[
      s + m = 1.
    \]

    \begin{enumerate}[(a)]
      \item (20 points) Without any intervention, what proportion of music
            $m^*$ does Lucas choose?

      \item (30 points) The community's utility function is
            \[
              U(s,m) = 1 - (s^2 + 3m^2),
            \]
            reflecting that the community values silence more than music.
            \begin{enumerate}[(i)]
              \item Calculate the proportion of music $m_c$ the community
                    would like Lucas to play. Explain why it differs from
                    $m^*$.
              \item Is $(s_c, m_c)$ Pareto efficient? Identify \emph{all}
                    values of $m$ that constitute Pareto efficient allocations.
            \end{enumerate}

      \item (20 points) Find the socially optimal allocation $(s_s, m_s)$
            that maximises the unweighted aggregate welfare
            $W = u(s,m) + U(s,m)$.

      \item (30 points) The community decides Lucas's music proportion should
            be lower. Knowing his utility is measured in dollars, elaborate
            and describe \textbf{two mechanisms} that could induce him to
            decrease his choice of $m$ and explain how each alters his
            incentives.
    \end{enumerate}
  },
  choices   = {},
  answer    = {},
  solution  = {%
    \textbf{(a)} Substituting $s = 1-m$:
    \[
      u = 1 - (1-m)^2 - m^2 = -2m^2 + 2m.
    \]
    First-order condition: $-4m + 2 = 0 \Rightarrow m^* = \tfrac{1}{2}$.

    \textbf{(b.i)}
    \[
      U = 1 - (1-m)^2 - 3m^2 = -4m^2 + 2m.
    \]
    F.O.C.: $m_c = \tfrac{1}{4}$. Lucas plays more music than the community
    prefers because he does not internalise the negative noise externality
    he imposes ($U = u - 2m^2$, where $-2m^2$ is the community's noise
    disutility).

    \textbf{(b.ii)} $(s_c, m_c) = (\tfrac{3}{4}, \tfrac{1}{4})$ is Pareto
    efficient. All $m \in [\tfrac{1}{4}, \tfrac{1}{2}]$ are Pareto
    efficient: for $m < \tfrac{1}{4}$ both utilities increase if $m$
    rises; for $m > \tfrac{1}{2}$ both decrease if $m$ rises; in between,
    one utility rises and the other falls.

    \textbf{(c)}
    \[
      W = u + U = -6m^2 + 4m.
    \]
    F.O.C.: $m_s = \tfrac{1}{3}$, $s_s = \tfrac{2}{3}$.

    \textbf{(d)} Two possible mechanisms:
    \begin{itemize}
      \item \textbf{Pigouvian tax}: charge Lucas a tax of $t$ per unit of
            music played (e.g.\ on $m - \tfrac{1}{4}$ when $m \ge
            \tfrac{1}{4}$). This raises the marginal private cost of music,
            internalising the noise externality.
      \item \textbf{Silence subsidy}: pay Lucas $\sigma$ per unit of silence
            (e.g.\ on $s > \tfrac{3}{4}$). This increases the marginal
            private benefit of silence, making music relatively less
            attractive.
    \end{itemize}
  },
}
